
\begin{frame}{UDE – Uninformed}
        \par\textbf{Premise}: \textbf{Universal Differential Equations}. 
        \par\textbf{Model Layers}: Can we exploit past Architecture?.
        \begin{figure}
            \centering
            \includegraphics[width=0.4\textwidth]{stock2.png}
        \end{figure}
        
\end{frame}


\begin{frame} {UDE – Uninformed - Considerations}
    \begin{block} {Matrix Inversion}
        \par The equations of motion can be rewritten as:
        $$\ddot{\mathbf{q}} = \mathbf{M}(\mathbf{q})^{-1} [\mathbf{C}(\mathbf{q}, \dot{\mathbf{q}}) \dot{\mathbf{q}} + \mathbf{D} \dot{\mathbf{q}} + \mathbf{G}(\mathbf{q})]$$
    \end{block}
    \vspace{1em}
    \begin{itemize}
        \item Much time to converge: RK4 integration combined with deep backpropagation.
        \item Not constrained : no physical knowledge, not even good sense.
        \item \textbf{Initial Results}: Often \textbf{poor results} due to instability.
        \item \textbf{Ill-Posing}: Many \textbf{degrees of indeterminacy}.
    \end{itemize}
\end{frame}




\begin{frame}{Lagrangian Neural Networks (LNN) - Architecture}
    Euler-Lagrange equation can be \textbf{discretized} and expanded with chain rule

    $$ \frac{d}{dt} \nabla_{\dot{q}} \mathcal{L} - \nabla_{q} \mathcal{L} = 0 \rightarrow
    \bigl(\nabla_{\dot{q}} \nabla_{\dot{q}}^{\top} \mathcal{L}\bigr) \ddot{q} +
    \bigl(\nabla_{q} \nabla_{\dot{q}}^{\top} \mathcal{L}\bigr) \dot{q} = \nabla_{q} \mathcal{L} $$ 

    Since the hessian $\nabla_{\dot{q}} \nabla_{\dot{q}}^{\top} \mathcal{L}$, which is the \textbf{Inertia matrix}, is positive definite:

    $$ \ddot{q} = \bigl(\nabla_{\dot{q}} \nabla_{\dot{q}}^{\top} \mathcal{L}\bigr)^{-1} \biggl[ \nabla_{q} \mathcal{L} - \bigl(\nabla_{q} \nabla_{\dot{q}}^{\top} \mathcal{L}\bigr) \dot{q}  \biggr]$$ 

    $\mathcal{L}$ is the output of a black box neural network, while the Lagrangian Equation is computed through internal chains of gradients in the network

    \begin{figure}[T]
            \centering
            \includegraphics[width=0.3\linewidth]{stock.png} \\
            \vspace{-1em}
    \end{figure} 

\end{frame}


\begin{frame}{Lagrangian Neural Networks (LNN) - Results}
    \begin{columns}[T]
        \begin{column}{0.5\textwidth}
            \centering
            \textbf{Evolution testing}
            \begin{figure}[h]
                    \centering
                    \includegraphics[width=0.4\linewidth]{stock.png} \\
                    \vspace{-1em}
            \end{figure}    
        
        \end{column}

        \begin{column}{0.5\textwidth}
            \centering
            \textbf{Energy drift LNN vs baseline}
            \begin{figure}[T]
                    \centering
                    \includegraphics[width=0.4\linewidth]{stock.png} \\
                    \vspace{-1em}
            \end{figure} 
            
        \end{column}
    \end{columns}
\end{frame}


\begin{frame}{Lagrangian Neural Networks (LNN) - Considerations}
    \begin{itemize}
        \item Almost perfect \textbf{Energy} conservation
        \item positive definiteness of \textbf{Hessian} is not guaranteed being $\mathcal{L}$ black box, causing \textbf{instability}
        \item \textbf{Batch Size}: Very small $\rightarrow$ need for \textbf{noisy gradients} to escape shallow loss regions (as suggested in literature).
        \item \textbf{Speed Issues}: Hessian computation is expensive.
        \item \textbf{Anomaly}: Difficulty learning \textbf{non-conservative systems} or systems with \textbf{broken physics}.

        \begin{columns}[T]
            \begin{column}{0.5\textwidth}
                \centering
                \begin{figure}[h]
                        \centering
                        \includegraphics[width=0.4\linewidth]{stock.png} \\
                        \vspace{-1em}
                        \caption{loss descent - correct physics}
                \end{figure}    
            
            \end{column}

            \begin{column}{0.5\textwidth}
                \centering
                \begin{figure}[T]
                        \centering
                        \includegraphics[width=0.4\linewidth]{stock.png} \\
                        \vspace{-1em}
                        \caption{loss descent - broken physics}
                \end{figure} 
                
            \end{column}
        \end{columns}


    \end{itemize}
\end{frame}

